\section{Enforcing Non-Negativity}\label{sec:nonneg}

Section~\ref{sec:reduction} reduces each candidate free set to an SPD solve, leaving
$x \geq 0$ as the sole remaining constraint. Two strategies enforce it. The first, an
active-set loop, wraps the CG solver unchanged and certifies global optimality through the
complementarity system~\eqref{eq:lcp}; the second, projected gradient, handles the constraint
inside every step but forgoes the Krylov rate.

\subsection{Active-set / block-principal-pivoting loop}\label{ssec:activeset}

The active-set method maintains a free set $\mathcal{F}$, solves~\eqref{eq:innerspd} (or its
equality-augmented form) on $\mathcal{F}$ by CG, and adjusts $\mathcal{F}$ toward
feasibility. In the \emph{primal step} any free variable that returned negative is pushed to
the bound and dropped from $\mathcal{F}$. Once all free variables are non-negative, the
\emph{dual step} inspects the reduced gradient $s = Ax - b$ (minus $B^\top\lambda$ in the
equality-augmented case, with $\lambda$ from Section~\ref{sec:reduction}) at every bound
variable: a bound index $i$ with $s_i < 0$ would lower $f$ if released, so it is re-admitted
to $\mathcal{F}$. The loop terminates when no variable violates either test, i.e.\ when
$x \geq 0$, $s \geq 0$ and complementary slackness hold to tolerance -- exactly the KKT
system~\eqref{eq:lcp}, which for the strictly convex $f$ certifies the unique global minimiser.

Toggling an index between $\mathcal{F}$ and $\mathcal{B}$ is a \emph{principal pivot} on
$\mathrm{LCP}(A,\,{-b})$; dropping or adding several indices at once is a \emph{block}
principal pivot~\cite{judice1994,kim2011}. Because $A \succ 0$ is a $P$-matrix, every
principal submatrix is invertible and each pivot is well defined~\cite{murty1988}. Block
pivots are fast but, like any all-at-once exchange, can cycle: a batch drop can over-shoot the
support and a later batch add restore it, revisiting a free set already seen.
Algorithm~\ref{alg:elim} therefore runs the batch exchange as a \emph{fast path} guarded by an
anti-cycling counter. Let $N$ be the current number of primal and dual violators and $\bar N$
the fewest seen so far. A batch step is taken whenever it makes progress ($N < \bar N$) or a
patience budget $p$ is unspent; otherwise the algorithm falls back to a single
\emph{Bland-style} pivot on the lowest-indexed violator. The least-index single pivot is
Murty's principal pivoting method, which cannot cycle~\cite{murty1988}; it is what the
finite-termination guarantee rests on, the batch path serving only to reach the optimum faster
when it does not stall. This is the anti-cycling construction of J\'udice and Pires for block
principal pivoting~\cite{judice1994}.

\begin{algorithm}[ht]
\caption{Active-set outer loop (wraps the CG inner solver $f$)}\label{alg:elim}
\begin{algorithmic}[1]
\Require Inner solver $f$ returning $(x_{\mathcal{F}}, k_{\mathrm{step}})$ for free set $\mathcal{F}$;\;
         SPD $A$ (or its mat-vec) and $b$;\; tolerance $\varepsilon$;\;
         patience $p_{\max} \in \mathbb{N}$ \textit{(default $3$)}
\Ensure Minimiser $x \in \mathbb{R}^n$ of~\eqref{eq:bqp};\; outer steps $s$;\; total inner iterations $k$
\State $\mathcal{F} \leftarrow \{1,\ldots,n\}$;\; $s \leftarrow 0$;\; $k \leftarrow 0$;\;
       $\bar N \leftarrow n+1$;\; $p \leftarrow p_{\max}$
\Loop
  \State $(x_{\mathcal{F}},\, k_{\mathrm{step}}) \leftarrow f(\mathcal{F})$;\;
         $s \leftarrow s + 1$;\; $k \leftarrow k + k_{\mathrm{step}}$
    \hfill\textit{($k_{\mathrm{step}}$: CG iterations; $1$ for a direct inner solve)}
  \State $x_i \leftarrow (x_{\mathcal{F}})_i$ for $i \in \mathcal{F}$;\; $x_i \leftarrow 0$ for $i \notin \mathcal{F}$
  \State $s_i \leftarrow (Ax)_i - b_i$
    \hfill\textit{(reduced gradient; subtract $(B^\top\lambda)_i$ in the equality-augmented case)}
  \State $\mathcal{D} \leftarrow \{i \in \mathcal{F} : x_i < -\varepsilon\}$;\;
         $\mathcal{V} \leftarrow \{i \notin \mathcal{F} : s_i < -\varepsilon\}$
    \hfill\textit{(primal and dual violators)}
  \State $\mathcal{H} \leftarrow \mathcal{D} \cup \mathcal{V}$;\; $N \leftarrow |\mathcal{H}|$
  \If{$N = 0$}
    \State \textbf{break} \hfill\textit{(KKT~\eqref{eq:lcp} satisfied; globally optimal)}
  \ElsIf{$N < \bar N$ \textbf{ or } $p > 0$}
    \hfill\textit{(fast path: progress, or patience remains)}
    \State \textbf{if} $N < \bar N$ \textbf{ then } $\bar N \leftarrow N$;\; $p \leftarrow p_{\max}$ \textbf{ else } $p \leftarrow p - 1$
    \State $\mathcal{F} \leftarrow (\mathcal{F} \setminus \mathcal{D}) \cup \mathcal{V}$
      \hfill\textit{(batch exchange: drop all $\mathcal{D}$, add all $\mathcal{V}$)}
  \Else
    \hfill\textit{(anti-cycling fallback: $p$ exhausted)}
    \State $i^\star \leftarrow \min\{i : i \in \mathcal{H}\}$
      \hfill\textit{(Bland least-index rule)}
    \State \textbf{if} $i^\star \in \mathcal{D}$ \textbf{ then } $\mathcal{F} \leftarrow \mathcal{F} \setminus \{i^\star\}$
           \textbf{ else } $\mathcal{F} \leftarrow \mathcal{F} \cup \{i^\star\}$
      \hfill\textit{(single principal pivot)}
  \EndIf
\EndLoop
\State \Return $x$, $s$, $k$
\end{algorithmic}
\end{algorithm}

\begin{theorem}[Finite convergence of Algorithm~\ref{alg:elim}]\label{thm:termination}
Let $f(x) = \tfrac{1}{2}x^\top A x - b^\top x$ with $A \succ 0$. Then for any patience budget
$p_{\max} \in \mathbb{N}$, Algorithm~\ref{alg:elim} terminates in finitely many outer
iterations and returns the unique global minimiser of $\min_{x \geq 0} f(x)$. No
non-degeneracy assumption is required. The same holds for the equality-augmented
problem~\eqref{eq:eqp} with the reduced gradient of line~5 adjusted by $B^\top\lambda$,
provided $B_{\mathcal{F}}$ retains full row rank on every free set the loop visits (automatic
for $p=1$; generic once $|\mathcal{F}| \geq p$).
\end{theorem}

\begin{proof}
\emph{Step 1 (each inner solve is well posed).}
Because $A \succ 0$, every principal submatrix $A_{\mathcal{F}}$ is SPD, so $A$ is a
$P$-matrix~\cite{murty1988}, and the system $A_{\mathcal{F}}x_{\mathcal{F}} = b_{\mathcal{F}}$
has a unique solution to which CG converges (Proposition~\ref{prop:cg}). In the
equality-augmented case the free-set solve is the saddle system~\eqref{eq:saddle}; with
$A_{\mathcal{F}} \succ 0$ and $B_{\mathcal{F}}$ of full row rank its Schur complement $S$ is
SPD, so~\eqref{eq:saddle} has the unique solution recovered by the $p+1$ SPD solves of
Section~\ref{sec:reduction}.

\emph{Step 2 (termination implies global optimality).}
Free-set stationarity gives $s_i = 0$ for $i \in \mathcal{F}$. The loop exits only when
$N = 0$: (i) all free $x_i \geq 0$ and (ii) all bound $s_i \geq 0$. These are exactly
$x \geq 0$, $s \geq 0$, $x_i s_i = 0$ -- the KKT system~\eqref{eq:lcp}. Strict convexity of
$f$ makes them sufficient for the unique global minimiser; it therefore suffices to show the
loop cannot run forever.

\emph{Step 3 (the fallback cannot cycle).}
Call a maximal block of consecutive iterates taking the \textbf{else} branch (the single Bland
pivot) a \emph{fallback run}. Within a run $\bar N$ is constant, and each iterate performs one
principal pivot on the lowest-indexed violator. This is precisely Murty's least-index
principal pivoting method for the $P$-matrix LCP of Step~2, which generates a sequence of
complementary bases (here, free sets) in which no basis recurs and which reaches the solution
in finitely many pivots~\cite{murty1988}. A fallback run is an initial segment of that
sequence started from the free set left by the preceding step, so each run is finite,
wherever it starts.

\emph{Step 4 (only finitely many batch steps and runs).}
The counter satisfies $0 \le N \le n$, and $\bar N$ is non-increasing, strictly decreasing
exactly when the progress branch ($N < \bar N$) fires; hence that branch fires at most $n+1$
times. Each remaining fast-path step has $N \ge \bar N$ and decrements $p$, and $p$ resets only
on a progress step, so at most $p_{\max}$ such steps occur between progress events. The number
of fallback runs is likewise bounded: a run ends at termination ($N=0$) or at a progress step,
and there are at most $n+1$ progress steps. Thus the algorithm takes at most
$(n+1) + (n+2)\,p_{\max}$ fast-path steps and at most $n+2$ fallback runs.

\emph{Step 5 (finite termination).}
By Step~3 each of the at most $n+2$ fallback runs is finite, and by Step~4 the fast-path steps
are finite in number; their sum is finite, so the loop reaches $N=0$ after finitely many outer
iterations, and Step~2 certifies the returned $x$ as the unique global minimiser. A crude
ceiling is the number of complementary bases, $2^n$.
\end{proof}

The $2^n$ ceiling establishes \emph{finiteness} only; it is not an efficiency estimate. The
guarantee is unconditional: the least-index fallback removes the non-degeneracy hypothesis that
earlier active-set arguments need, and it is the fallback -- not the batch exchange -- that the
proof rests on. In practice the fallback is rarely reached, because a batch step fails to
reduce $N$ only when a dropped index is immediately re-added (or vice versa), which forces
$s_i = 0$ exactly for some toggled $i$ -- a codimension-one, hence non-generic,
condition~\cite{judice1994}; block principal pivoting therefore converges in a handful of outer
steps on typical data, with the fast path alone certifying optimality.

\begin{remark}[Inexact inner solves]\label{rem:exact}
Theorem~\ref{thm:termination} assumes each inner call returns the exact minimiser of
$f_{\mathcal{F}}$. CG is run to a relative residual tolerance $\varepsilon$, so the returned
iterate is inexact; because the dual-feasibility test of line~5 uses the same $\varepsilon$, an
iterate accurate to $\varepsilon$ is accepted. The finite-termination argument holds exactly in
exact arithmetic and to working accuracy in finite precision.
\end{remark}

\subsection{Projected-gradient comparator}\label{ssec:proj}

The second strategy enforces the constraint at every step. Projected gradient iterates
\[
  x_{k+1} = \Pi_C\!\left(x_k - \tau\,\nabla f(x_k)\right),
  \qquad \nabla f(x) = A x - b,
\]
where $C$ is the feasible set. For~\eqref{eq:bqp} it is the non-negative orthant, so
$\Pi_C(y) = \max(y, 0)$ componentwise; for the single normalisation $\mathbf{1}^\top x=\beta$,
$x\geq0$ it is the simplex slice $\Delta_\beta$, projected onto in $O(n\log n)$ by a single
sort-and-threshold pass~\cite{duchi2008,condat2016}. For a general equality system $Bx=c$,
however, the projection onto $\mathcal{P} = \{x\geq0: Bx=c\}$ is itself a non-negative
quadratic program with no closed form, so projected gradient loses its per-step simplicity --
one reason the active-set loop, whose cost rises only by the $p$ extra right-hand sides of
Section~\ref{sec:reduction}, is the method of choice once $p>1$. The step size is $\tau = 1/L$ with
$L = \lambda_{\max}(A)$, estimated once by power iteration on the same (matrix-free) operator.
There is no outer loop and no linear solve: each step is one gradient evaluation -- two products
with $M$ when $A = M^\top M$ -- plus a projection.

\begin{proposition}[Projected-gradient convergence~{\cite{nesterov2004}}]\label{prop:proj}
  Let $A \succ 0$, so $f$ is $\mu$-strongly convex and $L$-smooth with
  $\mu = \lambda_{\min}(A)$ and $L = \lambda_{\max}(A)$, and let
  $x^\star = \arg\min_{x \in C} f(x)$. Projected gradient with $\tau = 1/L$ satisfies
  \[
    \|x_k - x^\star\|^2 \;\leq\; \left(1 - \tfrac{1}{\kappa(A)}\right)^{\!k}\|x_0 - x^\star\|^2,
    \qquad k \geq 0.
  \]
\end{proposition}
\begin{proof}[Proof sketch]
$x^\star = \Pi_C(x^\star - \tau\nabla f(x^\star))$ is a fixed point; non-expansiveness of
$\Pi_C$ and co-coercivity of the $L$-smooth gradient give
$\|x_{k+1}-x^\star\|^2 \le \|x_k-x^\star\|^2 - \tau(2-\tau L)\langle \nabla f(x_k)-\nabla f(x^\star),\,x_k-x^\star\rangle$.
Setting $\tau = 1/L$ makes $2-\tau L = 1$, and $\mu$-strong convexity bounds the inner product
below by $\mu\|x_k-x^\star\|^2$, so $\|x_{k+1}-x^\star\|^2 \le (1-\mu/L)\|x_k-x^\star\|^2$.
\end{proof}

The iteration count scales as $O(\kappa)$, against $O(\sqrt\kappa)$ for the CG inner loop
(Proposition~\ref{prop:cg}). Both methods pay the same per-step cost -- one operator
application -- so the $\sqrt\kappa$ gap in iteration counts is a runtime gap on
ill-conditioned problems. Nesterov acceleration would lift the projected-gradient rate to
$O(\sqrt\kappa)$, but CG attains that rate \emph{optimally} within the Krylov subspace: each
iterate minimises $f$ over $\mathcal{K}_k(A_{\mathcal{F}}, b_{\mathcal{F}})$, a strictly richer
search space than two-point momentum, and typically with a smaller constant. Plain projected
gradient is therefore the honest first-order comparator, isolating the Krylov advantage; a
FISTA variant~\cite{beck2009} narrows the constant but not the asymptotic gap. The trade-off
is otherwise in the other direction: projected gradient is loop-free and needs no
complementarity certificate, so it is the simpler method when $\kappa$ is small or moderate
accuracy suffices.

\subsection{Warm-starting}\label{ssec:warm}

For a parametric family of problems -- e.g.\ a sequence $b(\theta_1), b(\theta_2), \dots$ of
right-hand sides, or a swept normalisation $\beta$ -- consecutive minimisers usually differ in
only a few active constraints. Passing the previous problem's $(\mathcal{F}, x)$ pair as the
initial state of the next solve reduces both the outer count (the free set needs fewer primal
and dual adjustments) and the inner count (CG starts from a small initial residual, its guess
being the previous solution restricted to the new free set and renormalised). A direct inner
solver profits only from the warm free set, having no analogue of an initial guess; the
matrix-free CG inner solver profits from both, which is where the parametric sweeps of the
companion papers~\cite{schmelzer2026matrixfree,schmelzer2026rmt} obtain their largest speed-ups.
